% Opening material before first \section

The main finding across all of the (see appendix B) above cases (specifically, where the total replicative ``rate-density'' remains fixed between comparrisons), both in the continuous time and discrete time cases, is that the mean of the process at the end of the specified period is identical.  How it gets there varies, along with what happens after. Also, the variance and distribution seems to vary.  But the point remains -- early replicative enphasis given constant overall "potential replicative density", gives no advantage over a specified period, to the expected overall growth. \ 
Of course, this assumption of eqivalent "replicative density" does not apply to the case of comparrison between complete etracellular growth when set against the alternative of an initial "replicative boost" in macrophages. 











\subsection*{Two common exceptions to this rule}

\subsubsection{Non-conservation}
Non-finite supply of replicative means. Basically when the key assumption is removed. Perhaps a higher replication rate is accessible at larger population sizes.



\subsubsection{Monopoly effect}
Matthew effect






In appendix B, there is a summary of some work witch mathematically demonstrates a property of branching. That articifially front weighted replicative advantage --- when assuming an overall constant "replicative potental" --- does not give any advantage to a process with regards to its expected growth. In more intuitive terms, say a bacterial population had access to a finite supply of food on an agar plate, if it was to use more of it towards the beginning, and less later, there would not be any advantage to its overall expected growth. This seems relatively obvious, but I just wanted to check, and the maths checks out. At least in the mean. When it comes to the variance and distribution of outcomes, there is some nuance,  but let us leave this aside. \par 

What of the case of plague finding early replicative in macrophages? The condition of constant "replicative potential" does not apply when comparing with the case of less fecund extracellular replication.
